% GENERATED FILE. Edit canonical lessons, then run npm run generate:books.
% canonical-source-hash: fnv1a64:65f240abb84da1bd
% canonical-lessons: TA-C155-tallu, TA-C155-ilu, TA-C155-potu, TA-C155-kattu2, TA-C155-tuvai

\chapter{To push, To pull, To put, To tie, To wash}
\label{ch:ta-to-push-to-pull-to-put-to-tie-to-wash}
\hlchaptermodality{155}

\begin{chapteropening}
\textbf{By the end of this chapter:} \emph{I can say to push, to pull, to put, to tie and to wash.}

Complete the last lesson of chapter 155: i can say to push, to pull, to put, to tie and to wash.
\end{chapteropening}

\section[taḷḷu]{\ta{தள்ளு} (taḷḷu) — to push}

\label{lesson:TA-C155-tallu}



\begin{quote}
Before the new one: say the Tamil for to count, then the Tamil for to shout.
\end{quote}

\subsection*{You'll want to know: \ta{தள்ளு}}
\textbf{\ta{தள்ளு}} — \emph{taḷḷu} — \textquotedblleft{}to push\textquotedblright{}.

\begin{grammarlens}[title={What you've built}]
One more word for this chapter.
\end{grammarlens}

\subsection*{Your turn}
\begin{itemize}
\raggedright
  \item \emph{Say it:} \emph{taḷḷu}
  \item \emph{Say it:} \emph{taḷḷu}, once more
  \item \emph{Say it:} say \emph{taḷḷu}
  \item \emph{Recall:} say the Tamil for to count, then the Tamil for to shout, then say \emph{taḷḷu} again
\end{itemize}

\begin{tcolorbox}[breakable,colback=teal!4,colframe=teal!35!black,title={Before you move on}]
What does \ta{தள்ளு} mean? (\textquotedblleft{}To push\textquotedblright{}.) Say it once more.
\end{tcolorbox}
\section[iḻu]{\ta{இழு} (iḻu) — to pull}

\label{lesson:TA-C155-ilu}



\begin{quote}
Before the new one: say the Tamil for to shout, then the Tamil for to push.
\end{quote}

\subsection*{You'll want to know: \ta{இழு}}
\textbf{\ta{இழு}} — \emph{iḻu} — \textquotedblleft{}to pull\textquotedblright{}.

\begin{grammarlens}[title={What you've built}]
One more word for this chapter.
\end{grammarlens}

\subsection*{Your turn}
\begin{itemize}
\raggedright
  \item \emph{Say it:} \emph{iḻu}
  \item \emph{Say it:} \emph{iḻu}, once more
  \item \emph{Say it:} say \emph{iḻu}
  \item \emph{Recall:} say the Tamil for to shout, then the Tamil for to push, then say \emph{iḻu} again
\end{itemize}

\begin{tcolorbox}[breakable,colback=teal!4,colframe=teal!35!black,title={Before you move on}]
What does \ta{இழு} mean? (\textquotedblleft{}To pull\textquotedblright{}.) Say it once more.
\end{tcolorbox}
\section[pōṭu]{\ta{போடு} (pōṭu) — to put, to drop}

\label{lesson:TA-C155-potu}



\begin{quote}
Before the new one: say the Tamil for to push, then the Tamil for to pull.
\end{quote}

\subsection*{You'll want to know: \ta{போடு}}
\textbf{\ta{போடு}} — \emph{pōṭu} — \textquotedblleft{}to put, to drop\textquotedblright{}.

\begin{grammarlens}[title={What you've built}]
One more word for this chapter.
\end{grammarlens}

\subsection*{Your turn}
\begin{itemize}
\raggedright
  \item \emph{Say it:} \emph{pōṭu}
  \item \emph{Say it:} \emph{pōṭu}, once more
  \item \emph{Say it:} say \emph{pōṭu}
  \item \emph{Recall:} say the Tamil for to push, then the Tamil for to pull, then say \emph{pōṭu} again
\end{itemize}

\begin{tcolorbox}[breakable,colback=teal!4,colframe=teal!35!black,title={Before you move on}]
What does \ta{போடு} mean? (\textquotedblleft{}To put, to drop\textquotedblright{}.) Say it once more.
\end{tcolorbox}
\section[kaṭṭu]{\ta{கட்டு} (kaṭṭu) — to tie, to build}

\label{lesson:TA-C155-kattu2}



\begin{quote}
Before the new one: say the Tamil for to pull, then the Tamil for to put.
\end{quote}

\subsection*{You'll want to know: \ta{கட்டு}}
\textbf{\ta{கட்டு}} — \emph{kaṭṭu} — \textquotedblleft{}to tie, to build\textquotedblright{}.

\begin{grammarlens}[title={What you've built}]
One more word for this chapter.
\end{grammarlens}

\subsection*{Your turn}
\begin{itemize}
\raggedright
  \item \emph{Say it:} \emph{kaṭṭu}
  \item \emph{Say it:} \emph{kaṭṭu}, once more
  \item \emph{Say it:} say \emph{kaṭṭu}
  \item \emph{Recall:} say the Tamil for to pull, then the Tamil for to put, then say \emph{kaṭṭu} again
\end{itemize}

\begin{tcolorbox}[breakable,colback=teal!4,colframe=teal!35!black,title={Before you move on}]
What does \ta{கட்டு} mean? (\textquotedblleft{}To tie, to build\textquotedblright{}.) Say it once more.
\end{tcolorbox}
\section[tuvai]{\ta{துவை} (tuvai) — to wash (clothes)}

\label{lesson:TA-C155-tuvai}



\begin{quote}
Before the new one: say the Tamil for to put, then the Tamil for to tie.
\end{quote}

\subsection*{You'll want to know: \ta{துவை}}
\textbf{\ta{துவை}} — \emph{tuvai} — \textquotedblleft{}to wash (clothes)\textquotedblright{}.

\begin{grammarlens}[title={What you've built}]
One more word for this chapter.
\end{grammarlens}

\subsection*{Your turn}
\begin{itemize}
\raggedright
  \item \emph{Say it:} \emph{tuvai}
  \item \emph{Say it:} \emph{tuvai}, once more
  \item \emph{Say it:} say \emph{tuvai}
  \item \emph{Recall:} say the Tamil for to put, then the Tamil for to tie, then say \emph{tuvai} again
\end{itemize}

\begin{tcolorbox}[breakable,colback=teal!4,colframe=teal!35!black,title={Before you move on}]
What does \ta{துவை} mean? (\textquotedblleft{}To wash (clothes)\textquotedblright{}.) Say it once more.
\end{tcolorbox}
